\summarychapter{33}{The Heart and Its Reasons}{Pascal and the Pietists}

\begin{inbrief}
Even in the Age of Reason the thirst of the soul could not be ignored. The Catholic Blaise Pascal and
the mostly Lutheran Pietists belong together because both held that the heart is essential to knowing
God. The chapter tells of the Jansenist protest against the easy morals taught by French Jesuits, of
Pascal's conversion and writings, and of Spener, Francke and Zinzendorf, whose Pietism became the
fountain of modern revivals.
\question{How, in the Age of Reason, did Pascal and the Pietists sustain the life of the soul?}
\end{inbrief}

A few days after Pascal's death a servant noticed a bulge in the great scientist's jacket. In the
lining was a parchment in his hand, with a scribbled paper copy: ``The year of grace 1654. Monday,
November 23rd, from about half past ten in the evening until about half past twelve. God of Abraham,
God of Isaac, God of Jacob, not of the philosophers and scholars. Certainty, certainty, feeling, joy,
peace\dots{} Renunciation, total and sweet.'' It was the record of two hours in the presence of God.
For eight years he had carried it hidden, sewing it into each new coat, his private memorial of the
coming of grace. With faith under attack, old labels meant less. Pascal's critique showed that the
thinkers of the Enlightenment could not defend their narrow, technical idea of reason without large
assumptions, leaps of faith of their own; his witness rested not on a self-sufficient reason but on
the deeper wisdom that links reason to the heart.

\ssection{Appeals to ``Cheap Grace''}
The Catholic church of the seventeenth century had its heroes in Europe and overseas, but in France
it guarded its alliances with power. The higher clergy shared with nobles and king the nation's
revenues and the splendor of the court. The Jesuits, teachers of the wealthy and confessors of the
powerful, became expert psychologists. They learned to soften the thunder of Sinai and to ease
remorse by casuistry, the science of right and wrong for the Christian conscience. A confessor had to
know what is mortal sin, venial sin or no sin, when the letter of the law may yield to its spirit,
and when a Christian might lie, steal or kill for a higher purpose. Some casuists were strict; the
Jesuits were very lenient. They made so many allowances for human weakness that earnest people
protested against ``cheap grace,'' forgiveness without contrition.

\begin{sidebox}{Mental Reservation and Probabilism}
Wide mental reservation lets an ambiguous statement deceive the hearer. When Roman soldiers were
chasing Athanasius, his friend told them he was ``not far away,'' though the bishop was hidden in the
boat, and the soldiers hurried on. Strict mental reservation lets the speaker mislead by a
qualification kept in his own mind, for example that a lie told for a very noble reason need not be
admitted. Jesuits argued that such an answer is neither a lie nor a sin. They also proposed
probabilism, the theory that an act is justified if there is even a remote probability that it is
morally acceptable.
\end{sidebox}

The strongest opposition came from Jansenism. Cornelius Jansen (1585--1638), a Dutchman, adopted
Augustine's views of sin and grace at the University of Louvain. He believed the best defense of
Catholicism against Calvinism was a return to Augustine and a rigorous moral code for the clergy. He
fought the Jesuits as professor of Scripture at Louvain and then as bishop of Ypres, and left a great
treatise, \emph{Augustinus}, published in 1640, which became the platform of Jansenism in France. It
taught that God chose before the foundation of the world those who would be saved, and that good
works cannot earn salvation without grace, because the will is not free and human nature is corrupt
beyond self-rescue. Only God's grace, available through Christ's death, saves. The church had never
repudiated these doctrines but had let them slip into the background, and after Trent the Jesuit
defense of free will so magnified human works that, as Jansen saw it, it cheapened Christ's
redemption. The Jesuits made too much of reason and too little of trusting faith.

\ssection{Jansenism at Port-Royal}
Jansen's friend Jean DuVergier, abbot of Saint-Cyran, brought the message to France in 1633 as
confessor of Port-Royal, a Cistercian convent sixteen miles from Paris. Its young abbess, Jacqueline
Arnauld, led her nuns in wholehearted and humble devotion, and its reputation drew other women and a
group of devout laymen, the Solitaries, who sought holiness in the marshy valley around it. Her
brother Antoine Arnauld, of the Sorbonne, the theological faculty of Paris, took over the leadership
after Jansen's death. In 1643 he attacked, without naming the Jesuits, the idea that frequent
confession can make up for frequent sinning. The Jesuits persuaded the pope that Jansenism was
Calvinism in Catholic dress, and in 1653 he condemned five propositions said to be taken from
\emph{Augustinus}. Arnauld attacked again in two \emph{Letters to a Duke and Peer}, and the Sorbonne
moved to expel him. He needed a defender and turned to Blaise Pascal (1623--1662).

Pascal had known little but trouble. His mother died when he was three, and his father Stephen
educated the three children, Gilberte, Blaise and Jacqueline, and took the boy to meetings of the
Academy of Science. Before he was twenty-seven Blaise had won the admiration of the mathematicians of
Paris, invented a calculating machine for his father, a tax collector, and discovered the basic
principles of atmospheric and hydraulic pressure. He met Jansenism through an accident. In January
1646 Stephen fell on icy ground while hurrying to stop a duel and dislocated a hip. The two
physicians who treated him were Jansenists, and they won the son. They taught that bodily suffering
illustrates a religious truth: man by himself is helpless and miserable. Blaise, who had seldom
passed a day without pain, felt the force of it. He also learned from them how deeply the Bible
speaks to the human condition and began to study it as he had studied air pressure. In 1651 his
father died; his sister Jacqueline entered Port-Royal, and he was left alone in Paris. He took a
richly furnished house, servants and a coach and four and pursued polite society as he had pursued
geometry. After a year he felt only ``great disgust with the world'' and thought himself abandoned by
God.

\ssection{Certainty, Joy, Peace}
In that darkness he turned to the seventeenth chapter of John, and the Fire came, as his memorial
records. Late in 1654 he joined his sister in the community of Port-Royal, where Arnauld found him.
Pascal answered with eighteen \emph{Provincial Letters}, exposing Jesuit theology and practice with
eloquence and sarcastic wit. They were instant best sellers and made an obscure convent the center of
public attention. The pope condemned them; all educated France read them, and went on reading them
for two centuries. After the last letter, in March 1657, he planned a book on the evidences of
Christianity but never finished it. He fell violently ill in June 1662 and died on August 19, aged
thirty-nine.

Eight years later friends published his notes on faith and reason as the \emph{Pens\'ees}
(\emph{Thoughts}). Here he cuts across doctrine to the moral problem of man. Nature, he insisted, is
no sure clue to God: he sees ``too much to deny and too little to assure.'' In this ambiguous
universe man is the greatest mystery, part angel and part beast, like the Chimera of Greek myth: ``a
novelty, a monster, a chaos, a contradiction, a prodigy! The glory and refuse of the universe.''
Reason is no sure guide either. Trusting it alone we would doubt everything except pain and death,
yet the heart tells us that life cannot be meaningless. God and the meaning of life must be felt by
the heart: ``The heart has its reasons which reason does not know.'' Three centuries later people
still take their bearings from him. After his death the church and Louis XIV together drove Jansenism
from France. Port-Royal was destroyed, and the movement took refuge in Holland.

France was not the only land in need of heart. Any religion that becomes the habit of a majority
tends to go flat, and so it was in much of Lutheran Germany.

\begin{sidebox}{Protestant Scholasticism}
In the hundred years after Luther and Calvin the Protestant movement changed. The early call to
entrust oneself to God's transforming mercy was eclipsed by doctrinal discussion that seemed abstract
and lifeless. Luther read the Bible for its life-giving witness to Jesus; the Protestant Scholastics
built long arguments on the logic of Aristotle, which Luther had condemned. When they appealed to
Scripture they lifted out statements as isolated bits of logical data and lost the great flow of
God's story. Scripture came to be seen as a self-sufficient body of material, a source, no longer a
witness made alive by the Spirit. Theology was shaped by the schools more than by the church, and
sermons sounded like lectures. What Luther and Calvin had said mattered less than what they would have
said on speculative topics. All this did not distinguish Protestants from unbelievers or even from
Catholics but from one another, and confessions and creeds divided the church more and more.
\end{sidebox}

After the creative Reformation came the cautious period called Protestant Scholasticism or
confessionalism. Luther's robust justifying faith became, among his followers, a mental exercise:
formal assent to doctrines set out by scholars in place of surrender to the mercy of God in Christ.
The Christian life became membership in the state church, and the marks of a good Christian were
regular attendance and reception of the sacraments from orthodox ministers. Pietism arose against
this hardening, as Jansenism had against cheap grace. It had two aims. First, personal faith. The
Pietists left behind the dreams of Catholic Christendom and Puritan commonwealths and began with the
individual, so that for the first time the conversion of baptized Christians, and not only of pagans,
came to the fore; the essence of faith is a personal experience of God's grace in the heart. Second,
they wanted to move the center of Christian life from the state churches into which people were born
to close fellowships of those with living faith, from which renewed laymen would carry the Word to
all classes. Three men mark its history.

\ssection{The Soul's Return to Life}
Philip Jacob Spener (1635--1705) grew up on the German mystic Johann Arndt and the English Puritans,
and at the University of Strasbourg learned to understand justification by faith as a spiritual
rebirth. After three years as a pastor there he was called to Frankfurt, where conditions shocked
him. He left the prescribed texts, preached from the whole Bible and called for repentance and
discipleship. Nothing remarkable happened until 1669, when his sermons on the Sermon on the Mount
brought sudden conversions and changed families. Twice a week he gathered a small company in his
house for Scripture reading and religious conversation. People called them in scorn ``gatherings of
the pious,'' and Pietism was born. His introduction to an edition of Arndt's sermons, \emph{Pious
Desires}, recommended Bible study groups for spiritual growth, a strenuous and rather ascetic life,
more care for the character of theological students, and simpler, more spiritual preaching. A cell of
experiential Christians in each congregation would be leaven for the whole, and only those born again
through conscious conversion were fit for the ministry. In 1686 he became court preacher at Dresden,
where his plain preaching brought trouble, and in 1692 he moved to Berlin at the invitation of the
elector of Brandenburg. The same year, with wisdom and humility, he had August Hermann Francke
(1663--1727) called to the new University of Halle. Francke soon led the movement, though Spener wrote
and preached until his death in February 1705.

Francke came from Leipzig, where his spiritual emphasis had caused tension with colleagues. His zeal
went back to a conversion after two years of doubt. Writing a sermon on John 20:31, he fell to his
knees and begged the God he did not yet know or believe in to save him, if he existed; sadness left
him in ``a flood of joy,'' and afterward he found it easy to live righteously and gladly. At Halle
he made the university the hub of Pietist work: a school for the poor, an orphanage, a hospital built
on the site of a tavern he bought, a Latin school for gifted boys, homes for widows and for unmarried
women, a medical dispensary, a book depot, a printing house and a Bible house. When Frederick IV of
Denmark wanted to found one of the earliest Protestant missions, in India in 1705, he found his first
missionaries among Francke's students.

\ssection{The Carpenter and the Count}
The later phase belongs to Count Nikolaus von Zinzendorf (1700--1760), an ardent, emotional man for
whom true Christianity is a simple, childlike faith in the blood of Jesus, and who sang of the soul's
relation to Christ in vivid, almost erotic images. He was the grandson of a nobleman who had left
Austria for his faith and was raised by a devout grandmother who admired Halle. He studied law at
Wittenberg for three years, but his heart was in ministry. His opportunity came through the remnant
of the old Hussite movement, the Bohemian Brethren (\emph{Unitas Fratrum}), who had flourished at the
Reformation, been nearly crushed in the Thirty Years' War and were still persecuted. Stirred by
revival under a Moravian carpenter, Christian David, they sought refuge in Protestant lands and found
it on Zinzendorf's estate, where in 1722 David founded Herrnhut, ``The Lord's Watch.'' Devout people
of many backgrounds flocked there. The Moravians intended a town of Christians only, separate from
the world, a real communion of saints---a free, social monasticism without celibacy.

From 1727 Zinzendorf was its guiding spirit, and ten years later he was ordained in the reorganized
Moravian Church, or United Brethren. His impulses were always missionary, and the Moravians became
the first large-scale Protestant missionary force in history. At the coronation of Christian VI in
Copenhagen in 1731 he met a black man from the Danish West Indies and was struck by the needs of the
enslaved. In 1732 Leonhard Dober and David Nitschmann left Herrnhut for St.~Thomas, and missions
followed quickly in Greenland, Lapland, among the Indians of Georgia, on the Guinea Coast, among the
Hottentots of South Africa, in Guiana, Ceylon and Algeria. Zinzendorf's religion ran to sentiment,
yet some of his hymns, such as ``Jesus, still lead on,'' are widely sung, and few have shown such
devotion to Christ: ``I have only one passion. It is He, none but He.''

\ssection{View from a Rearview Mirror}
Pietism gave much to Christianity worldwide. It turned eighteenth-century churches from controversy
to the care of souls, made preaching and pastoral visiting central to the Protestant ministry,
enriched Christian music and showed how much a revived church needs a spiritual laity. Its greatest
legacy may be the small group and the devotional reading of Scripture. Personal copies of the Bible
were more available than ever before, personal engagement with it became the focus of discipleship,
and groups would dwell on a word from Scripture reinforced by a well-chosen hymn. Under all this lay
its dominant theme, regeneration, meaning not a doctrine but the indispensable experience of every
Christian, the true fulfillment of the Reformation: in the heart doctrine becomes reality. Its
intensely personal language made Christianity a drama of the soul, the heart a battlefield of good
and evil. In this sense Pietism is the fountain of all modern revivals, and evangelical Christianity
is unthinkable without it.

Evangelicals inherited two of its traits. First, emotion played so large a part that reason was
endangered. Feelings were left to carry the meaning of faith, so Pietism had little to say about
God's place in nature and history and offered little resistance to secularism; Evangelicals often
share the weakness. Second, Pietism took the institutional church for granted and did not attack it,
but it moved what is essential, the new birth and the spiritual life, from the state churches to
fellowship groups and voluntary associations. In the United States, where no state church prevailed,
Evangelicals allied this with the denominational idea and multiplied movements devoted to evangelism
or some aspect of the Christian life. Pietism lives on in the Moravian church and in evangelical
Christianity at large, the spiritual descendants of John Wesley and George Whitefield.

\begin{keyfacts}
\keypeople{Blaise Pascal; Cornelius Jansen; the abbot of Saint-Cyran; Jacqueline Arnauld; Antoine
Arnauld; Louis XIV; Johann Arndt; Philip Jacob Spener; August Hermann Francke; Count von Zinzendorf;
Christian David; Leonhard Dober and David Nitschmann.}
\keydates{%
  \kd{1640}{Jansen's \emph{Augustinus} published}
  \kd{1653}{The pope condemns five Jansenist propositions}
  \kd{1654}{Pascal's night of fire (November 23)}
  \kd{1662}{Death of Pascal}
  \kd{1669}{Spener's preaching at Frankfurt brings the first Pietist gatherings}
  \kd{1692}{Francke called to the University of Halle}
  \kd{1705}{Death of Spener; Halle Pietists go as missionaries to India}
  \kd{1722}{Herrnhut founded on Zinzendorf's estate}
  \kd{1732}{First Moravian missionaries sail for St.~Thomas}}
\keyterms{Casuistry; mental reservation; probabilism; ``cheap grace''; Jansenism; Port-Royal;
\emph{Provincial Letters}; \emph{Pens\'ees}; Protestant Scholasticism (confessionalism); Pietism;
\emph{Pious Desires}; regeneration; Herrnhut; Moravians.}
\end{keyfacts}
